\documentclass{article}
\usepackage{mathtools} 
\usepackage{fontspec}
\usepackage[UTF8]{ctex}
\usepackage{amsthm}
\usepackage{mdframed}
\usepackage{xcolor}
\usepackage{amssymb}
\usepackage{amsmath}


% 定义新的带灰色背景的说明环境 zremark
\newmdtheoremenv[
  backgroundcolor=gray!10,
  % 边框与背景一致，边框线会消失
  linecolor=gray!10
]{zremark}{说明}

% 通用矩阵命令: \flexmatrix{矩阵名}{元素符号}{行数}{列数}
\newcommand{\flexmatrix}[4]{
  \[
  #1 = \begin{pmatrix}
    #2_{11}     & #2_{12}     & \cdots & #2_{1#4}   \\
    #2_{21}     & #2_{22}     & \cdots & #2_{2#4}   \\
    \vdots      & \vdots      & \ddots & \vdots     \\
    #2_{#31}    & #2_{#32}    & \cdots & #2_{#3#4}
  \end{pmatrix}
  \]
}

% 简化版命令（默认矩阵名为A，元素符号为a）: \quickmatrix{行数}{列数}
\newcommand{\quickmatrix}[2]{\flexmatrix{A}{a}{#1}{#2}}

\begin{document}
\title{注释 1.3}
\author{张志聪}
\maketitle

\begin{zremark}
  例2 另一种解法。
\end{zremark}

参考《陶哲轩实分析》

命题前面的处理是相同的，证明$c^2 = a$的方式不同：

利用反证法，我们来证明$c^2 < a$和$c^2 > a$都会产生矛盾。

(i) 假设$c^2 < a$。设$0 < \epsilon < 1$是一个较小的正数，
由于$a \leq c \leq 1$且$\epsilon^2 \leq \epsilon$，所以
\begin{align*}
  (c + \epsilon)^2 & = c^2 + 2c\epsilon + \epsilon^2 \\
                   & = c^2 + 2c\epsilon + \epsilon   \\
                   & \leq c^2 + (2c + 1) \epsilon
\end{align*}
因为$c^2 < a$，所以通过适当的选取$\epsilon$
（利用实数的阿基米德性质），可得得到
\begin{align*}
  c^2 + (2c + 1) \epsilon < a
\end{align*}
即
\begin{align*}
  (c + \epsilon)^2 < a
\end{align*}
于是可得
\begin{align*}
  c + \epsilon \in E
\end{align*}
存在矛盾。

(ii) 假设$c^2 > a$，思路一致，
\begin{align*}
  (c - \epsilon)^2 & = c^2 - 2c\epsilon + \epsilon^2  \\
                   & \geq c^2 - 2c\epsilon 
\end{align*}
因为$c^2 > a$，所以通过适当的选取$\epsilon$，可得得到
\begin{align*}
  c^2 - 2c\epsilon > a
\end{align*}
即
\begin{align*}
  (c - \epsilon)^2 > a 
\end{align*}
所以，$\forall x \in E$，有
\begin{align*}
  c - \epsilon \geq x
\end{align*}
（因为，如果$c - \epsilon < x$，那么$(c - \epsilon)^2 < x^2 <a$，这是一个矛盾）。

所以，$c - \epsilon$是$E$的一个上界，存在矛盾。


\end{document}